\section{Challenges and Opportunities}
\label{Sec7:Challenges and Opportunities}
Studies of LLM-MA frameworks and applications are 
advancing rapidly, giving rise to numerous challenges and opportunities.   We identified several critical challenges and potential areas for future study.  

%\subsection{Agents-Environment Interface: Advancing into the Multi-Modal Environment}
\subsection{Advancing into  Multi-Modal Environment}

Most previous work on LLM-MA has been focused on text-based environments, excelling in processing and generating text. However, there is a notable lack in multi-modal settings, where agents would interact with and interpret data from multiple sensory inputs and generate multiple outputs such as images, audio, video, and physical actions. Integrating LLMs into multi-modal environments presents additional challenges, such as processing diverse data types and enabling agents to understand each other and respond to more than just textual information. 

%\subsection{Agents Communication: Hallucination}
\subsection{Addressing Hallucination}
% What is hallucination
The hallucination problem is a significant challenge in LLMs and single LLM-based Agent systems. It refers to the phenomenon where the model generates text that is factually incorrect~\cite{huang2023survey}. However, this problem takes on an added layer of complexity in a multi-agent setting. In such scenarios, one agent's hallucination can have a cascading effect. This is due to the interconnected nature of multi-agent systems, where misinformation from one agent can be accepted and further propagated by others in the network. 
Therefore, detecting and mitigating hallucinations in LLM-MA is not just a crucial task but also presents a unique set of challenges. It involves not only correcting inaccuracies at the level of individual agents but also managing the flow of information between agents to prevent the spread of these inaccuracies throughout the system.

% Single Agent => Multi-Agents
% Potential challenges and solutions
%Different from the single LLM-based agent, in the Multi-Agents scenario, hallucinations can have cascading effects due to the complexity of managing multiple agents. This means that false information from one agent can propagate through the system as other agents might accept and act upon it. The hallucination may be expanded during the agent communication. Hence, detecting and mitigating the hallucination in the LLM-based Multi-Agents is a promising and challenging task.

%\subsection{Agents Capabilities Acquisition}
\subsection{Acquiring Collective Intelligence}
%As we discussed in Section~\ref{Sec3:Taxonomy}, for the feedback of the Multi-Agents Learning, % off-policy Learn from Dataset, Model
% on-policy
In traditional multi-agent systems,  agents often use reinforcement learning to learn from offline training datasets.
However, LLM-MA systems mainly learn from instant feedback, such as interactions with the environment or humans, as we discussed in Section~\ref{Sec3:Taxonomy}.
%different from Multi-Agents reinforcement learning which can learn from offline training dataset, current LLM-based Multi-Agents learning is limited to learn from instant feedback such as environment or human. 
This learning style requires a reliable interactive environment and it would be tricky to design such an interactive environment for many tasks, limiting the scalability of LLM-MA systems. Moreover, the prevailing approaches in current research involve employing Memory and Self-Evolution techniques to adjust agents based on feedback. While effective for individual agents, these methods do not fully capitalize on the potential collective intelligence of the agent network. They adjust agents in isolation, overlooking the synergistic effects that can emerge from coordinated multi-agent interactions. Hence,  jointly adjusting multiple agents and achieving optimal collective intelligence is still a critical challenge for LLM-MA.

% Effective Knowledge sharing; partial knowledge sharing
% More effective synergy 
% effective and efficient 
% trade-off between different partial knowledge sharing 


% \subsection{Scaling law of the Multi-Agents}
\subsection{Scaling Up LLM-MA Systems}
\label{scale_up}
% challenges: computational complexity; scaling law; Coordination

LLM-MA systems are composed of a number of individual LLM-based agents, posing a significant challenge of scalability regarding the number of agents. 
%single LLM-based agents, which raises a critical challenge named the scaling of the Multi-Agents.
% computational complexity
From the computational complexity perspective, each LLM-based agent, typically built on large language models like GPT-4, demands substantial computational power and memory. Scaling up the number of these agents in an LLM-MA system significantly increases  resource requirements. In scenarios with limited computational resource, it would be challenging to develop these LLM-MA systems.
%Achieving the performance of Multi-Agents in a limited computational resource would be a major challenge.

% Coordination, Communication, Learning and Adaptation
Additionally, as the number of agents in an LLM-MA system increases, additional complexities and research opportunities emerge, particularly in areas like efficient agent coordination, communication, and understanding the scaling laws of multi-agents. For instance, with more LLM-based agents, the intricacy of ensuring effective coordination and communication rises significantly. 
{As highlighted in ~\cite{multiorches2}, designing advanced Agents Orchestration methodologies is increasingly important. These methodologies aim to optimize agents workflows, task assignments tailored to different agents, and communication patterns across agents such as communication constraints between agents. Effective Agents Orchestration facilitates harmonious operation among agents, minimizing conflicts and redundancies.}
Additionally, exploring and defining the scaling laws that govern the behavior and efficiency of multi-agent systems as they grow larger remains an important area of research. These aspects highlight the need for innovative solutions to optimize LLM-MA systems, making them both effective and resource-efficient.

%\subsection{Multi-Agents Datasets and Benchmarks}
\subsection{Evaluation and Benchmarks}
% Agent Profiling Role-playing-Following ability
% Communication quality
We have summarized the datasets and benchmarks currently available for LLM-MA in  Table~\ref{db}. This is a starting point, and far from  being comprehensive. We identify two significant challenges in evaluating  LLM-MA systems and benchmarking their performance against each other. 
Firstly, as discussed in ~\cite{xu2023magic}, much of the existing research focuses on evaluating individual agents' understanding and reasoning within narrowly defined scenarios. This focus tends to overlook the broader and more complex emergent behaviors that are integral to multi-agent systems. 
Secondly, there is a notable shortfall in the development of comprehensive benchmarks across several research domains, such as Science Team for Experiment Operations, Economic analysis, and Disease propagation simulation. This gap presents an obstacle to accurately assessing and benchmarking the full capabilities of LLM-MA systems in these varied and crucial fields.
%we conclude two challenges for MultiAgents datasets and benchmarks. The first challenge is as discussed in ~\cite{xu2023magic}, most previous work focuses on evaluating the ultimate understanding and reasoning capability of agents in specific settings, overlooking the emergent capabilities in multi-agent systems. Another challenge is that in some research areas such as the Science Team for Experiment Operations, Economy, Disease propagation simulation, etc, there is still a lack of comprehensive benchmarks to evaluate the capabilities of LLM-MA.


\subsection{Applications and Beyond}
% \paragraph{Application:}
The potential of LLM-MA systems extends far beyond their current applications, holding great promise for advanced computational problem-solving in fields such as finance, education, healthcare,  environmental science, urban planning and so on. As we have discussed, LLM-MA systems possess the capability to tackle complex problems and simulate various aspects of the real world. While the current role-playing capabilities of LLMs may have limitations, ongoing advancements in LLM technology suggest a bright future. It is anticipated to have more sophisticated methodologies, applications, datasets, and benchmarks tailored for diverse research fields.
Furthermore, there are opportunities to explore LLM-MA systems from various theoretical perspectives, such as Cognitive Science~\cite{sumers2023cognitive}, Symbolic Artificial Intelligence, Cybernetics, Complex Systems, and Collective Intelligence.  Such a multi-faceted approach could  contribute to a more comprehensive understanding and innovative applications in this rapidly evolving field.

% \paragraph{Multi-Agents Systems Design and Behaviors Understanding from Diverse Perspectives:} O
% Cognitive 
% Complex Systems
